\chapter{Liver Abscess}
\index{Liver Abscess}

\section*{Definition and Overview}
A liver abscess is a pus-filled mass within the liver parenchyma, most commonly caused by bacteria (pyogenic liver abscess) or parasites (amoebic liver abscess, typically \textit{Entamoeba histolytica}), and rarely by fungi. It represents a serious clinical condition that can lead to rapid systemic sepsis, tissue destruction, and death if untreated. Early recognition, diagnostic imaging, and prompt initiation of targeted antimicrobial therapy and drainage are key components of management.

\section*{Epidemiology}
Liver abscesses are uncommon but clinically significant. Pyogenic liver abscesses (PLAs) are more prevalent in developed nations like some countries, with an annual incidence of about 2 to 4 cases per 100,000 population. They occur more frequently in males and in patients aged 50 years and older. In contrast, amoebic liver abscesses (ALAs) are more common in tropical and subtropical regions with poor sanitation and are typically diagnosed in travellers returning from endemic areas or immigrants. Biliary tract disease is the most common predisposing factor for PLA, followed by diabetes mellitus, which increases the risk of PLA by two- to threefold.

\section*{Aetiology and Pathophysiology}
Pyogenic liver abscesses develop via several routes: biliary tract obstruction, portal vein bacteraemia, systemic circulation, direct extension from adjacent infections, or trauma. Amoebic abscesses result from bowel wall invasion and portal migration.
\begin{itemize}
    \item \textbf{Biliary tract disease:} Ascending cholangitis, choledocholithiasis, or biliary malignancies lead to bile stasis and bacterial proliferation, which ascends into the hepatic parenchyma.
    \item \textbf{Portal vein bacteraemia:} Intra-abdominal infections such as appendicitis or diverticulitis seed bacteria into the portal venous drainage, leading to hepatic microabscesses that coalesce.
    \item \textbf{Causative organisms (Pyogenic):} Common pathogens include \textit{Escherichia coli}, \textit{Klebsiella pneumoniae} (particularly virulent strains causing metastatic infections), \textit{Streptococcus} species, and anaerobic organisms.
    \item \textbf{Parasitic invasion (Amoebic):} \textit{Entamoeba histolytica} trophozoites colonise the colon, penetrate the mucosa, enter the portal circulation, and reach the liver, where they cause liquefactive necrosis.
    \item \textbf{Haematogenous spread:} Systemic bacteraemia via the hepatic artery from distant sites (e.g. endocarditis, osteomyelitis, or intravenous drug use).
\end{itemize}

\section*{Clinical Features}
The presentation of liver abscess is often non-specific, characterised by constitutional symptoms and localised abdominal signs.
\begin{itemize}
    \item \textbf{Fever and constitutional symptoms:} High-grade fever, spiking chills, night sweats, anorexia, weight loss, and malaise are present in the majority of patients.
    \item \textbf{Abdominal pain:} Dull, aching pain localised to the right upper quadrant (RUQ), which may radiate to the right shoulder due to diaphragmatic irritation.
    \item \textbf{Hepatomegaly and RUQ tenderness:} Tenderness upon palpation of the liver, and pain elicited by percussion over the lower right rib cage.
    \item \textbf{Jaundice:} Yellowing of the skin and sclera, which is particularly common in pyogenic abscesses secondary to biliary tract disease.
    \item \textbf{Pulmonary symptoms:} Cough, dyspnoea, or right-sided pleuritic chest pain due to reactive pleural effusion or an elevated hemidiaphragm.
\end{itemize}

\section*{Differential Diagnosis}
Liver abscesses must be differentiated from other inflammatory, infectious, and neoplastic conditions of the hepatobiliary system and abdomen.
\begin{itemize}
    \item \textbf{Acute cholecystitis and cholangitis:} Inflammatory conditions of the gallbladder or bile ducts that present with RUQ pain and fever but lack a parenchymal mass.
    \item \textbf{Primary hepatic malignancy or metastasis:} Hepatocellular carcinoma or metastatic lesions that may undergo central necrosis, mimicking an abscess on imaging.
    \item \textbf{Hepatic cysts:} Benign simple cysts, hydatid cysts (caused by \textit{Echinococcus granulosus}), or polycystic liver disease, which are typically non-inflammatory unless superinfected.
    \item \textbf{Pneumonia or empyema:} Right-sided lower lobe pulmonary infections that can cause referred RUQ pain and diaphragmatic irritation.
    \item \textbf{Amoebic vs. pyogenic abscess:} Differentiating the two is critical as management differs, with amoebic cases responding to antiparasitic therapy alone.
\end{itemize}

\section*{When to Seek Medical Care}
Patients with unexplained high fever, persistent right-sided abdominal pain, or yellowing of the skin and eyes should seek prompt medical evaluation.

\textbf{Contact your local emergency services for:}
\begin{itemize}
    \item High fever accompanied by confusion, severe lethargy, or rapid breathing
    \item Severe, intolerable abdominal pain or sudden worsening of pain
    \item Signs of shock, such as cold, clammy skin, rapid heart rate, or fainting
    \item Difficulty breathing or chest pain
\end{itemize}

\section*{Diagnosis and Investigations}
Diagnosis requires a combination of laboratory markers, microbiological cultures, and diagnostic imaging.
\begin{itemize}
    \item \textbf{Diagnostic imaging:} Transabdominal ultrasound is the initial imaging modality of choice; however, contrast-enhanced CT of the abdomen is more sensitive and characterises the abscess number, size, and location.
    \item \textbf{Microbiological cultures:} Blood cultures should be obtained in all suspected cases, along with culture of the abscess fluid obtained via aspiration.
    \item \textbf{Serological testing:} Enzyme-linked immunosorbent assay (ELISA) for anti-amoebic antibodies is highly sensitive and specific for diagnosing amoebic liver abscess.
    \item \textbf{Inflammatory markers:} Marked elevation of white cell count (leukocytosis), C-reactive protein (CRP), and erythrocyte sedimentation rate (ESR).
    \item \textbf{Liver function tests:} Elevated alkaline phosphatase (ALP) is the most common abnormality, along with mild elevations in transaminases and bilirubin.
\end{itemize}

\section*{Management}
Treatment of pyogenic liver abscess requires antibiotic therapy combined with drainage, while amoebic liver abscess is primarily managed pharmacologically.
\begin{itemize}
    \item \textbf{Percutaneous drainage:} Ultrasound- or CT-guided catheter drainage or needle aspiration is the gold standard for pyogenic abscesses larger than 3--5 cm.
    \item \textbf{Empirical antibiotic therapy:} Intravenous broad-spectrum antibiotics (e.g. ceftriaxone plus metronidazole, or piperacillin/tazobactam) started immediately after cultures are drawn.
    \item \textbf{Directed antibiotic therapy:} Tailoring antibiotics based on culture results, continuing for 4 to 6 weeks (typically 1--2 weeks intravenously followed by oral therapy).
    \item \textbf{Antiparasitic therapy:} Metronidazole or tinidazole is the primary treatment for amoebic liver abscess, followed by a luminal amoebicide (e.g. paromomycin) to eradicate intestinal colonisation.
    \item \textbf{Surgical intervention:} Laparoscopic or open surgical drainage is reserved for ruptured abscesses, multiloculated abscesses, or failure of percutaneous drainage.
\end{itemize}

\section*{Complications}
\begin{itemize}
    \item Rupture of the abscess into the peritoneal cavity, causing peritonitis and shock.
    \item Sepsis and septic shock, with systemic inflammatory response syndrome and multi-organ failure.
    \item Rupture through the diaphragm into the pleural space or pericardium, leading to empyema or cardiac tamponade.
    \item Metastatic infection, particularly endophthalmitis, meningitis, or brain abscess, most commonly associated with hypervirulent \textit{Klebsiella pneumoniae}.
    \item Portal vein thrombosis (pylephlebitis) due to spreading thrombophlebitis from the portal system.
\end{itemize}

\section*{Prognosis and Outcomes}
With modern diagnostic imaging, percutaneous drainage, and potent antimicrobial therapies, the mortality rate for pyogenic liver abscess has decreased from over 80\% in the mid-20th century to approximately 5\% to 10\%. Amoebic liver abscess has an even lower mortality rate (under 1\%) when treated early. Prognosis is less favourable in patients with multiple abscesses, malignant biliary obstruction, severe sepsis, or co-morbidities like advanced age and end-stage renal disease.

\section*{Prevention and Self-Care}
\begin{itemize}
    \item \textbf{Management of biliary disease:} Promptly treat gallstones and cholangitis to prevent ascending bacterial infections.
    \item \textbf{Travel hygiene:} Avoid drinking untreated water, eating unpeeled raw fruits and vegetables, or consuming unpasteurised dairy products in areas endemic for amoebiasis.
    \item \textbf{Glycaemic control:} Maintain tight blood glucose control in diabetic patients to reduce susceptibility to pyogenic infections.
    \item \textbf{Post-procedural monitoring:} Closely follow medical advice and report early signs of infection after biliary or abdominal surgeries.
    \item \textbf{Adherence to treatment:} Complete the full course of prescribed antibiotics to ensure complete resolution and prevent recurrence.
\end{itemize}

\section*{Sources and Further Reading}
The following reputable organisations provide further information on liver abscess and biliary health:
\begin{itemize}
    \item Gastroenterological Society of Australia. \textit{Information on liver abscess and biliary health.} https://www.gesa.org.au/
    \item Mayo Clinic. \textit{Liver abscess diagnosis and treatment.} https://www.mayoclinic.org/
    \item MSD Manual. \textit{Amoebic and pyogenic liver abscesses.} https://www.msdmanuals.com/
    \item NHS. \textit{Liver abscess causes, symptoms, and treatment.} https://www.nhs.uk/
    \item World Health Organization. \textit{Prevention and control of amoebiasis.} https://www.who.int/
\end{itemize}
